
\section{Equações de Hill/CW}
Como \eqref{eq:relativo_LVLHgeral} estava sendo expresso no quadro do \emph{target}, é obtido seu vetor posição:
\begin{equation}
\label{eq:posicao_LVLH_target}
\vec r_t = \begin{bmatrix} 0 \\ 0 \\ -r \end{bmatrix}
\end{equation}

E também o vetor de velocidade angular:
\begin{equation}
\label{eq:velocidadeAngular_LVLH_target}
\boldsymbol\omega = \begin{bmatrix} 0 \\ -\omega \\ 0
\end{bmatrix}
\end{equation}

Em termos da equação \eqref{eq:relativo_LVLHgeral}, é obtido:

\begin{equation}
\label{eq:velAngular_cross_s}
\boldsymbol\omega \times \boldsymbol s^* =
\begin{bmatrix}
    -\omega z \\
    0         \\
    \omega x
\end{bmatrix} 
\end{equation}

\begin{equation}
\label{eq:velAngular_cross_cross_s}
\boldsymbol\omega \times (\boldsymbol\omega \times \boldsymbol s^*) =
\begin{bmatrix}
    -\omega^2 x \\
    0           \\
    -\omega^2 z 
\end{bmatrix}
\end{equation}

\begin{equation}
\label{eq:omega_cross_dsdt}
\boldsymbol\omega \times \frac{d^* s^*}{dt} =
\begin{bmatrix}
    -\omega \dot z \\
    0                 \\
    \omega \dot x
\end{bmatrix}
\end{equation}

\begin{equation}
\label{eq:dOmega_cross_s}
\frac{d\boldsymbol\omega}{dt} \times \boldsymbol s^* =
\begin{bmatrix}
    -\dot\omega z \\
    0              \\
    \dot\omega x
\end{bmatrix}
\end{equation}

\begin{equation}
\label{eq:jacobiana_LVLH}
\mathbf Js^* =
\begin{bmatrix}
    1 & 0 & 0 \\
    0 & 1 & 0  \\
    0 & 0 & -2
\end{bmatrix} \boldsymbol s^* =
\begin{bmatrix}
    x \\
    y \\
    -2z
\end{bmatrix}
\end{equation}

No caso da órbita circular ($e=0$), a velocidade angular
\footnote{Os autores \cite{fehse2003,sidi1997spacecraft} definem a \emph{velocidade angular} (ou \emph{angular rate} ou até mesmo \emph{angular velocity}) como $\omega$, outros autores, como \cite{bate1971}, a definem como $\eta$} 
$\omega$ é constante e coincide com o movimento médio (\emph{mean motion}, $\eta$), definido por \cite{fehse2003}:

\begin{equation}
\label{eq:omega_constante}
\boldsymbol\omega^2 = \frac{\mu}{r_t^3}
\end{equation}

Ou podendo ser reescrita como:

\begin{equation}
\label{eq:omega_constante2}
\eta^2 = \frac{\mu}{r_t^3}
\end{equation}

Como $r_t=a$, a velocidade angular $\omega^{2}= \eta^2 = \mu/r_t^{3}$, portanto, neste caso em particular de órbita circular, $\omega=n$.

Então, $\dot{\omega}=0$, e ao inserir a equação \eqref{eq:omega_constante} nos termos da equação \eqref{eq:relativo_LVLHgeral}, é possível obter as equações linearizadas gerais do movimento relativo, conhecidas como Equações de Hill \cite{hill1878,fehse2003, wie2008space}.

\begin{equation}
    \label{eq:equacoes_HillX}
\ddot{x} - 2\omega \dot{z} = \frac{F_x}{m_c}
\end{equation}

\begin{equation}
    \label{eq:equacoes_HillY}
\ddot{y} + \omega^2 y = \frac{F_y}{m_c}
\end{equation}

\begin{equation}
    \label{eq:equacoes_HillZ}
\ddot{z} + 2\omega \dot{x} - 3 \omega^2 z = \frac{F_z}{m_c}
\end{equation}



\subsection{Mudança de Convenção de Quadro de Referência}

As equações de Hill apresentadas anteriormente foram formuladas na convenção do sistema de referência LVLH (\emph{Local Vertical Local Horizontal}), onde os eixos foram definidos da seguinte forma \cite{fehse2003}:

\begin{itemize}
    \item Eixo $x$: tangencial à órbita, na direção do movimento orbital (\emph{along-track}).
    \item Eixo $y$: normal à órbita, completando o sistema ortogonal  (\emph{cross-track}).
    \item Eixo $z$: radial, apontando para o centro da Terra  (\emph{radial}).
\end{itemize}

No entanto, para realizar a compatibilidade com documentos técnicos da NASA, será adotada a convenção de referência utilizada por ela, que define os eixos da seguinte forma \cite{Folta2022}:

\begin{itemize}
    \item Eixo $x$: radial, apontando do centro da Terra para o satélite (\emph{radial}).
    \item Eixo $y$: tangencial à órbita, na direção do movimento orbital  (\emph{along-track}).
    \item Eixo $z$: normal à órbita, completando o sistema ortogonal  (\emph{cross-track}).
\end{itemize}

Por conta dessa mudança de convenção, é necessário realizar uma rotação dos eixos e, consequentemente, uma reformulação das equações \eqref{eq:equacoes_HillX}, \eqref{eq:equacoes_HillY}, \eqref{eq:equacoes_HillZ}, conforme \cite{sidi1997spacecraft}:

\begin{equation}
    \label{eq:Hill_NASA_X}
    \ddot{x} - 2\omega \dot{y} - 3\omega^2 x = \frac{F_x}{m_c}
\end{equation}

\begin{equation}
    \label{eq:Hill_NASA_Y}
    \ddot{y} + 2\omega \dot{x} = \frac{F_y}{m_c}
\end{equation}

\begin{equation}
    \label{eq:Hill_NASA_Z}
    \ddot{z} + \omega^2 z = \frac{F_z}{m_c}
\end{equation}

As equações de Hill descrevem o movimento relativo de uma espaçonave em relação a outra em órbita próxima, considerando os efeitos de Coriolis, equação \eqref{eq:velAngular_cross_s} a força centrípeta - equação \eqref{eq:velAngular_cross_cross_s}. A convenção da NASA é utilizada em missões espaciais e facilita a integração com modelos dinâmicos e simulações padronizadas \cite{Folta2022}.


As equações diferenciais lineares com coeficientes que variam no tempo, conforme as equações \eqref{eq:equacoes_HillX}, \eqref{eq:equacoes_HillY} e \eqref{eq:equacoes_HillZ}, representam o modelo geral que é válido para uma trajetória relativa arbitrária \cite{debruijn2011,ortolano2017} entre o \emph{chaser} e o \emph{target}, considerando que este último se move sob a influência exclusiva de um campo gravitacional central, ou seja, pode ser considerado uma \emph{espaçonave não-cooperativa} \cite{mahendrakar2023}.

Com o objetivo de reduzir a ordem das matrizes, as equações de Hill fora reescritas na forma de espaço de estados. Foi adotada a separação do modelo em dois sistemas independentes: (i) uma para a dinâmica fora do plano, sendo a equação \eqref{eq:equacoes_HillY}, e (ii) outra para a dinâmica no plano, sendo as equações \eqref{eq:equacoes_HillX} e \eqref{eq:equacoes_HillZ}, as equações do movimento. A forma fundamental - ou geral - do espaço de estados pode ser expressa como \cite{fehse2003,starek2016,sherrill2014}:

\begin{equation}
    \dot{\mathbf{x}} = \mathbf{A}\,\mathbf{x} + \mathbf{B}\,\mathbf{u}
    \label{eq:state_space_general}
\end{equation}

Onde $\mathbf{A}$ é a matriz de transição, $\mathbf{x}$ é o vetor de estados, $\mathbf{B}$ é a matriz de entradas e $\mathbf{u}$ é o vetor de entradas, que neste caso é a força de controle aplicada no \emph{chaser}.

Para o primeiro caso dentro do plano (\textit{in-plane}), as equações provenientes de \eqref{eq:equacoes_HillX} e \eqref{eq:equacoes_HillZ} e considerando o vetor de estados $x = [x,z, \dot{x}, \dot{z}]^T$ podem ser escritas conforme \cite{fehse2003}:

% \begin{equation}
% \begin{bmatrix}
% \dot{x}(t) \\[0.3em]
% \dot{z}(t) \\[0.3em]
% \ddot{x}(t) \\[0.3em]
% \ddot{z}(t)
% \end{bmatrix}
% =
% \begin{bmatrix}
% 0 & 0 & 1 & 0 \\
% 0 & 0 & 0 & 1 \\
% 0 & 0 & 0 & 2\omega \\
% 0 & 3\omega^2 & -2\omega & 0
% \end{bmatrix}
% \begin{bmatrix}
% x(t) \\[0.3em]
% z(t) \\[0.3em]
% \dot{x}(t) \\[0.3em]
% \dot{z}(t)
% \end{bmatrix}
% +
% \begin{bmatrix}
% 0 & 0 \\
% 0 & 0 \\
% \frac{1}{m_c} & 0 \\
% 0 & \frac{1}{m_c}
% \end{bmatrix}
% \begin{bmatrix}
% F_x \\[0.3em]
% F_z
% \end{bmatrix}
% \label{eq:inplane}
% \end{equation}

\begin{equation}
    \label{eq:inplaneXY_NASA}
\begin{bmatrix}
\dot{x}(t) \\
\dot{y}(t) \\
\ddot{x}(t) \\
\ddot{y}(t)
\end{bmatrix}
=
\begin{bmatrix}
0 & 0 & 1 & 0 \\
0 & 0 & 0 & 1 \\
-3\omega^2 & 0 & 0 & -2\omega \\
0 & 0 & 2\omega & 0
\end{bmatrix}
\begin{bmatrix}
x(t) \\
y(t) \\
\dot{x}(t) \\
\dot{y}(t)
\end{bmatrix}
+
\begin{bmatrix}
0 & 0 \\
0 & 0 \\
\frac{1}{m_c} & 0 \\
0 & \frac{1}{m_c}
\end{bmatrix}
\begin{bmatrix}
F_x \\
F_y
\end{bmatrix}
\end{equation}

De forma similar, para a dinâmica fora do plano, (\textit{out-of-plane}), considerando o vetor de estados $x_0 = [z, \dot{z}]^T$, é obtido:

% \begin{equation}
%     \label{eq:outofplane}
% \begin{bmatrix}
% \dot{y}(t) \\[0.3em]
% \ddot{y}(t)
% \end{bmatrix}
% =
% \begin{bmatrix}
% 0 & 1 \\
% \omega^2 & 0
% \end{bmatrix}
% \begin{bmatrix}
% y(t) \\[0.3em]
% \dot{y}(t)
% \end{bmatrix}
% +
% \begin{bmatrix}
% 0 \\
% \frac{1}{m_c}
% \end{bmatrix}
% F_y
% \end{equation}

\begin{equation}
    \label{eq:outofplaneZ_NASA}
\begin{bmatrix}
\dot{z}(t) \\
\ddot{z}(t)
\end{bmatrix}
=
\begin{bmatrix}
0 & 1 \\
\omega^2 & 0
\end{bmatrix}
\begin{bmatrix}
z(t) \\
\dot{z}(t)
\end{bmatrix}
+
\begin{bmatrix}
0 \\
\frac{1}{m_c}
\end{bmatrix}
F_z
\end{equation}
